%% ============================================================
\section{\texorpdfstring{$\sigma^\mu$}{sigma} Completeness Relation}
%% ============================================================

\subsection{Main identity}

The $\sigma$ matrices form a complete basis for $2\times2$ complex matrices.
The completeness relation is (Srednicki eq.~38.1):
\begin{equation}
  \boxed{
    \sigma^\mu{}_{\alpha\dot\alpha}\,\sigma_\mu{}^{\beta\dot\beta}
    = -2\,\delta_\alpha{}^\beta\,\delta_{\dot\alpha}{}^{\dot\beta}
  }
  \tag{38.1}
\end{equation}
where $\sigma_\mu{}^{\beta\dot\beta} = g_{\mu\nu}\bar\sigma^{\nu\dot\beta\beta}$
(both spinor indices raised by $\varepsilon$).

\verified{Max error: $0.0e+00$.}

\medskip
The nonzero components of the LHS are:
$(\alpha,\dot\alpha,\beta,\dot\beta) \in
\{(1,1,1,1),(1,2,1,2),(2,1,2,1),(2,2,2,2)\}$, each equal to $-2$.

\subsection{Alternative form}

Equivalently (eq.~35.5):
\begin{equation}
  \varepsilon^{\alpha\beta}\varepsilon^{\dot\alpha\dot\beta}\,
  \sigma^\mu{}_{\alpha\dot\alpha}\,\sigma^\nu{}_{\beta\dot\beta}
  = -2\,g^{\mu\nu}
  \tag{35.5}
\end{equation}

\textbf{Cadabra2:} $\epsilon^{\alpha \beta} \epsilon^{\dal \dbe} \sigma^{\mu}\,_{\beta \dbe} = \bar\sigma^{\mu\dot\alpha\alpha}$

\verified{Max error vs $-2g^{\mu\nu}$: $0.0e+00$.}

%% ============================================================
\section{Trace Identities}
%% ============================================================

\subsection{Two-sigma trace}

The fundamental trace identity (Srednicki eq.~38.5):
\begin{equation}
  \boxed{
    \mathrm{Tr}[\sigma^\mu\bar\sigma^\nu]
    \equiv \sigma^\mu{}_{\alpha\dot\alpha}\,\bar\sigma^{\nu\dot\alpha\alpha}
    = -2\,g^{\mu\nu}
  }
  \tag{38.5}
\end{equation}

\textbf{Cadabra2 summand:} $\sigma^{\mu}\,_{\alpha \dal} \sigmabar^{\nu \dal \alpha}$

Explicit matrix (rows $\mu$, columns $\nu$):
\[
  \mathrm{Tr}[\sigma^\mu\bar\sigma^\nu] =
  \begin{pmatrix}
  +2.0 & +0.0 & +0.0 & +0.0 \\
  +0.0 & -2.0 & +0.0 & +0.0 \\
  +0.0 & +0.0 & -2.0 & +0.0 \\
  +0.0 & +0.0 & +0.0 & -2.0 \\
  \end{pmatrix}
  = -2g^{\mu\nu}
\]

\verified{$\mathrm{Tr}[\sigma^\mu\bar\sigma^\nu] = -2g^{\mu\nu}$: max error $0.0e+00$.}

\verified{$\mathrm{Tr}[\bar\sigma^\mu\sigma^\nu] = -2g^{\mu\nu}$: max error $0.0e+00$.}

\medskip
\textbf{Physical interpretation:}
$\sigma^\mu$ and $\bar\sigma^\nu$ are ``dual'' bases for $2\times 2$ matrices;
the trace gives their inner product $-2g^{\mu\nu}$.  Any $2\times2$ matrix $A$
expands as $A = -\frac{1}{2}\sum_\mu \bar\sigma^\mu\,\mathrm{Tr}[\sigma^\mu A]$.

\subsection{Four-sigma trace}

The four-sigma trace identity (eq.~38.6), with Srednicki's
mostly-plus metric $g = \mathrm{diag}(-1,+1,+1,+1)$:
\begin{equation}
  \boxed{
    \mathrm{Tr}[\sigma^\mu\bar\sigma^\nu\sigma^\rho\bar\sigma^\sigma]
    = 2\bigl(g^{\mu\nu}g^{\rho\sigma} - g^{\mu\rho}g^{\nu\sigma}
          + g^{\mu\sigma}g^{\nu\rho}\bigr)
      + 2i\,\varepsilon^{\mu\nu\rho\sigma}
  }
  \tag{38.6}
\end{equation}
\begin{equation}
  \mathrm{Tr}[\bar\sigma^\mu\sigma^\nu\bar\sigma^\rho\sigma^\sigma]
  = 2\bigl(g^{\mu\nu}g^{\rho\sigma} - g^{\mu\rho}g^{\nu\sigma}
        + g^{\mu\sigma}g^{\nu\rho}\bigr)
    - 2i\,\varepsilon^{\mu\nu\rho\sigma}
  \tag{38.6$'$}
\end{equation}
(The conjugate trace carries $-2i$ instead of $+2i$; the sign difference
is specific to the mostly-plus convention.)

\textbf{Sample numerical values} ($\varepsilon^{0123}=+1$):
\medskip
\begin{center}
\begin{tabular}{@{}llll@{}}
  \toprule
  $(\mu,\nu,\rho,\sigma)$ & Computed & Expected & \\
  \midrule
  $(0,1,2,3)$ & $+0.00++2.00i$ & $+0.00++2.00i$ & $\checkmark$ \\
  $(1,0,2,3)$ & $+0.00+-2.00i$ & $+0.00+-2.00i$ & $\checkmark$ \\
  $(0,1,3,2)$ & $+0.00+-2.00i$ & $+0.00+-2.00i$ & $\checkmark$ \\
  $(0,2,1,3)$ & $+0.00+-2.00i$ & $+0.00+-2.00i$ & $\checkmark$ \\
  $(1,2,3,0)$ & $+0.00+-2.00i$ & $+0.00+-2.00i$ & $\checkmark$ \\
  $(2,3,0,1)$ & $+0.00++2.00i$ & $+0.00++2.00i$ & $\checkmark$ \\
  \bottomrule
\end{tabular}
\end{center}
\verified{All 256 components verified. Max error: $0.0e+00$.}
\verified{Conjugate trace: max error $0.0e+00$.}

%% ============================================================
\section{Van der Waerden Index Gymnastics}
%% ============================================================

\subsection{Raising and lowering spinor indices}

The $\varepsilon$-metric raises and lowers spinor indices
(Srednicki eqs.~38.7--38.13):
\begin{align}
  \psi_\alpha &= \varepsilon_{\alpha\beta}\,\psi^\beta
  &\text{(lower undotted)},
  \tag{38.7}\\
  \psi^\alpha &= \varepsilon^{\alpha\beta}\,\psi_\beta
  &\text{(raise undotted)},
  \tag{38.8}\\
  \chi_{\dot\alpha} &= \varepsilon_{\dot\alpha\dot\beta}\,\chi^{\dot\beta}
  &\text{(lower dotted)},
  \tag{38.9}\\
  \chi^{\dot\alpha} &= \varepsilon^{\dot\alpha\dot\beta}\,\chi_{\dot\beta}
  &\text{(raise dotted)},
  \tag{38.10}
\end{align}
with the Srednicki convention
$\varepsilon_{12}=-1,\;\varepsilon^{12}=+1$
($\varepsilon_{21}=+1$, $\varepsilon^{21}=-1$).

Key consistency: $\varepsilon^{\alpha\beta}\varepsilon_{\beta\gamma} = \delta^\alpha_\gamma$.
\verified{$\varepsilon^{\alpha\beta}\varepsilon_{\beta\gamma}=\delta^\alpha_\gamma$: max error $0.0e+00$.}

\subsection{Raising both spinor indices of \texorpdfstring{$\sigma^\mu$}{sigma}}

The relation between $\sigma$ and $\bar\sigma$ via index raising (eq.~38.13):
\begin{equation}
  \bar\sigma^{\mu\dot\alpha\alpha}
  = \varepsilon^{\alpha\beta}\,\varepsilon^{\dot\alpha\dot\beta}\,
    \sigma^\mu{}_{\beta\dot\beta}
  \tag{38.13}
\end{equation}

\textbf{Cadabra2:} $\epsilon^{\alpha \beta} \epsilon^{\dal \dbe} \sigma^{\mu}\,_{\beta \dbe} = \sigmabar^{\mu \dal \alpha}$

\verified{All four $\mu$ values verified: max error $0.0e+00$.}

%% ============================================================
\section{Spinor Bilinears}
%% ============================================================

\subsection{Lorentz transformation properties}

The key bilinear forms and their Lorentz transformation properties:
\begin{align}
  \chi\psi &\equiv \chi^\alpha\psi_\alpha
  = \varepsilon^{\alpha\beta}\chi_\alpha\psi_\beta
  &\text{Lorentz scalar (left-handed)}
  \tag{38.14a}\\
  \chi^\dagger\psi^\dagger &\equiv \chi^\dagger_{\dot\alpha}\psi^{\dagger\dot\alpha}
  &\text{Lorentz scalar (right-handed)}
  \tag{38.14b}\\
  \chi^\dagger_{\dot\alpha}\,\bar\sigma^{\mu\dot\alpha\alpha}\,\psi_\alpha
  &
  &\text{Lorentz 4-vector}
  \tag{38.14c}
\end{align}

\textbf{Cadabra2 vector bilinear:} $\chidag_{\dal} \sigmabar^{\mu \dal \alpha} \psi_{\alpha}$

\subsection{Momentum matrix in spinor space}

The momentum $p^\mu$ maps to a $2\times2$ spinor matrix:
\begin{equation}
  p_{\alpha\dot\alpha} = p_\mu\,\sigma^\mu{}_{\alpha\dot\alpha}
  = -E\,I_2 + \mathbf{p}\cdot\boldsymbol\sigma,
\end{equation}
with det$(p_{\alpha\dot\alpha}) = -p^2/4$ (up to normalisation).
For massless particles: det$=0$, so $p_{\alpha\dot\alpha}$ factorises as
$p_{\alpha\dot\alpha} = \lambda_\alpha\tilde\lambda_{\dot\alpha}$ ---
the starting point for spinor-helicity methods.

\subsection{Gordon identity}

In 4-component Dirac notation, the vector bilinear decomposes as:
\begin{equation}
  \bar u(p')\,\gamma^\mu\,u(p)
  = \bar u(p')\!\left[\frac{p'^\mu+p^\mu}{2m}
    + \frac{i\sigma^{\mu\nu}(p'-p)_\nu}{2m}\right]\!u(p)
\end{equation}
In 2-component language, the vector part comes from
$\tilde u^{\dot\alpha s}\bar\sigma^{\mu\dot\alpha\alpha}u^s_\alpha$
and the tensor part from the generators $S^{\mu\nu}_L$.

%% ============================================================
\section{Fierz Identities}
%% ============================================================

\subsection{Schouten identity}

The $\varepsilon$-Fierz or Schouten identity:
\begin{equation}
  \boxed{
    \varepsilon_{\alpha\gamma}\,\varepsilon_{\beta\delta}
    = \varepsilon_{\alpha\beta}\,\varepsilon_{\gamma\delta}
    - \varepsilon_{\alpha\delta}\,\varepsilon_{\gamma\beta}
  }
  \tag{38.17}
\end{equation}

\textbf{Cadabra2:} $\epsilon_{\alpha \gamma} \epsilon_{\beta \delta} = \epsilon_{\alpha \beta} \epsilon_{\gamma \delta}\discretionary{}{}{}-\,\epsilon_{\alpha \delta} \epsilon_{\gamma \beta}$

\verified{Schouten identity: max error $0.0e+00$.}

\subsection{$\sigma$-matrix Fierz}

From the completeness relation (§1), the $\sigma$-matrix Fierz identity is:
\begin{equation}
  \bigl(\bar\sigma^\mu\bigr)^{\dot\alpha\alpha}\,\bigl(\sigma_\mu\bigr)_{\beta\dot\beta}
  = -2\,\delta^\alpha_\beta\,\delta^{\dot\alpha}_{\dot\beta}
  \tag{38.18}
\end{equation}

\subsection{Four-Fermi Fierz}

For four Grassmann-valued spinors:
\begin{equation}
  \bigl(\chi^\dagger_{\dot\alpha}\bar\sigma^{\mu\dot\alpha\alpha}\psi_\alpha\bigr)
  \bigl(\xi^\dagger_{\dot\beta}\bar\sigma_\mu{}^{\dot\beta\beta}\eta_\beta\bigr)
  = -2\,\bigl(\chi^\dagger_{\dot\alpha}\xi^{\dagger\dot\alpha}\bigr)
         \bigl(\eta^\beta\psi_\beta\bigr)
  \tag{38.21}
\end{equation}
This follows from the $\sigma$-Fierz plus Grassmann anticommutativity,
and is used to rewrite four-fermion interaction terms.

%% ============================================================
\section{Gordon Identity}
%% ============================================================

\subsection{Main identity}

The \textbf{Gordon identity} for on-shell Dirac spinors
(Srednicki Problem~38.3, eq.~38.22):
\begin{equation}
  \boxed{
    2m\,\bar u_{s'}(\mathbf{p}')\,\gamma^\mu\,u_s(\mathbf{p})
    = \bar u_{s'}(\mathbf{p}')\bigl[(p'+p)^\mu - 2i\,S^{\mu\nu}(p'-p)_\nu\bigr]
      u_s(\mathbf{p})
  }
  \tag{38.22}
\end{equation}
where $S^{\mu\nu} = \tfrac{i}{4}[\gamma^\mu,\gamma^\nu]$ is the Dirac spin tensor.
An analogous identity holds for the $v$ spinors with an overall sign flip.

\subsection{Derivation}

From the Clifford algebra and the definition of $S^{\mu\nu}$:
\begin{align}
  \gamma^\mu\slashed{p} &= -p^\mu - 2iS^{\mu\nu}p_\nu, \\
  \slashed{p}'\gamma^\mu &= -p^{\prime\mu} + 2iS^{\mu\nu}p'_\nu.
\end{align}
Summing:
\begin{equation}
  \gamma^\mu\slashed{p} + \slashed{p}'\gamma^\mu
  = -(p+p')^\mu - 2iS^{\mu\nu}(p_\nu - p'_\nu).
\end{equation}
Sandwiching with $\bar u_{s'}(\mathbf{p}')$ and $u_s(\mathbf{p})$, then using
$(\slashed{p}+m)u_s = 0$ and $\bar u_{s'}(\slashed{p}'+m) = 0$ gives the identity.

\subsection{Extended Gordon identity with \texorpdfstring{$\gamma^5$}{gamma5}}

Because $\{\gamma^5,\gamma^\mu\} = 0$ (eq.~36.46):
\begin{equation}
  \bar u_{s'}(\mathbf{p}')\bigl[(p'+p)^\mu - 2iS^{\mu\nu}(p'-p)_\nu\bigr]
  \gamma^5\,u_s(\mathbf{p}) = 0.
  \tag{38.41}
\end{equation}
The Gordon tensor is parity-even; $\gamma^5$ is parity-odd; their product vanishes.

\subsection{Physical content}

Decomposing the electromagnetic vertex $\gamma^\mu$:
\begin{equation}
  \bar u'\gamma^\mu u
  = \frac{(p+p')^\mu}{2m}\,\bar u'u
  + \frac{i\sigma^{\mu\nu}(p'-p)_\nu}{2m}\,\bar u'u
  \qquad
  (\sigma^{\mu\nu} = \tfrac{i}{2}[\gamma^\mu,\gamma^\nu])
\end{equation}
\begin{itemize}
  \item $(p+p')^\mu/(2m)$: charge-current term $\to$ form factor $F_1(q^2)$
  \item $i\sigma^{\mu\nu}q_\nu/(2m)$: magnetic-moment term $\to$ form factor $F_2(q^2)$, anomalous magnetic moment $g{-}2$
\end{itemize}

\verified{$\{\gamma^5,\gamma^\mu\}=0$ for all $\mu$: max error $<10^{-14}$.}
\verified{Gordon: $2m\bar u'\gamma^\mu u = \bar u'(2p^\mu)u$ at $p'=p$: max error $<10^{-14}$.}

%% ============================================================
\section{Helicity and Spin-Sum Completeness}
%% ============================================================

\subsection{Spin-sum completeness}

The Dirac spin-sum completeness relations (eqs.~38.28--38.29):
\begin{equation}
  \boxed{
    \sum_s u_s(p)\,\bar u_s(p) = -\slashed{p} + m\,,
    \qquad
    \sum_s v_s(p)\,\bar v_s(p) = -\slashed{p} - m\,.
  }
  \tag{38.28--29}
\end{equation}

\subsection{Bilinear identities}

For same momentum $\mathbf{p}$ (eqs.~38.21--38.22):
\begin{align}
  \bar u_{s'}(\mathbf{p})\,\gamma^\mu\,u_s(\mathbf{p})
    &= 2p^\mu\,\delta_{s's}, \tag{38.21a}\\
  \bar v_{s'}(\mathbf{p})\,\gamma^\mu\,v_s(\mathbf{p})
    &= 2p^\mu\,\delta_{s's}, \tag{38.21b}\\
  \bar u_{s'}(\mathbf{p})\,\gamma^0\,v_s(-\mathbf{p})
    &= 0, \tag{38.22a}\\
  \bar v_{s'}(\mathbf{p})\,\gamma^0\,u_s(-\mathbf{p})
    &= 0. \tag{38.22b}
\end{align}

\subsection{Helicity and chirality}

The \textbf{helicity} $h = \hat{\mathbf{J}}\cdot\hat{\mathbf{p}}$ measures spin
along the direction of motion.  For spin-$\tfrac{1}{2}$: $h = \pm\tfrac{1}{2}$.
In the massless limit, helicity coincides with chirality:
\begin{align}
  P_R\,u_+(\mathbf{p}) &= u_+(\mathbf{p}), & \text{(positive helicity = right-chiral)}\\
  P_L\,u_-(\mathbf{p}) &= u_-(\mathbf{p}), & \text{(negative helicity = left-chiral)}
\end{align}
where $P_{L,R} = \tfrac{1}{2}(1\mp\gamma^5)$.

The massless spin sums separate by helicity:
\begin{equation}
  u_+(\mathbf{p})\bar u_+(\mathbf{p}) = P_R(-\slashed{p}), \qquad
  u_-(\mathbf{p})\bar u_-(\mathbf{p}) = P_L(-\slashed{p}).
\end{equation}
Their sum gives the massless limit of eq.~(38.28): $\sum_s u_s\bar u_s \to -\slashed{p}$ as $m\to 0$.

\verified{Spin-sum $\sum_s u_s\bar u_s = -\slashed{p}+m$: max error $<10^{-14}$.}
\verified{Bilinear off-diagonal = 0; orthogonality $\bar u\gamma^0 v = 0$: verified.}
\verified{Massless chirality projectors: $P_Lu_- = u_-$, $P_Ru_+ = u_+$.}

%% ============================================================
\section{Numerical Verification Summary}
%% ============================================================

\begin{center}
\renewcommand{\arraystretch}{1.5}
\begin{tabular}{@{}lll@{}}
  \toprule
  Identity & Equation & Max error \\
  \midrule
  $\sigma^\mu{}_{\alpha\dot\alpha}\sigma_\mu{}^{\beta\dot\beta}=-2\delta_\alpha^\beta\delta_{\dot\alpha}^{\dot\beta}$
    & (38.1) & $0.0e+00$ \\
  $\varepsilon^{\alpha\beta}\varepsilon^{\dot\alpha\dot\beta}\sigma^\mu_{\alpha\dot\alpha}\sigma^\nu_{\beta\dot\beta}=-2g^{\mu\nu}$
    & (35.5) & $0.0e+00$ \\
  $\mathrm{Tr}[\sigma^\mu\bar\sigma^\nu]=-2g^{\mu\nu}$
    & (38.5) & $0.0e+00$ \\
  $\mathrm{Tr}[\bar\sigma^\mu\sigma^\nu]=-2g^{\mu\nu}$
    & (38.5) & $0.0e+00$ \\
  $\mathrm{Tr}[\sigma^\mu\bar\sigma^\nu\sigma^\rho\bar\sigma^\sigma]=2(\cdots)+2i\varepsilon$
    & (38.6) & $0.0e+00$ \\
  $\mathrm{Tr}[\bar\sigma^\mu\sigma^\nu\bar\sigma^\rho\sigma^\sigma]=2(\cdots)-2i\varepsilon$
    & (38.6$'$) & $0.0e+00$ \\
  $\varepsilon^{\alpha\beta}\varepsilon_{\beta\gamma}=\delta^\alpha_\gamma$
    & --- & $0.0e+00$ \\
  $\bar\sigma^{\mu\dot\alpha\alpha}=\varepsilon^{\alpha\beta}\varepsilon^{\dot\alpha\dot\beta}\sigma^\mu_{\beta\dot\beta}$
    & (38.13) & $0.0e+00$ \\
  Schouten identity & (38.17) & $0.0e+00$ \\
  Gordon: $2m\bar u'\gamma^\mu u = \bar u'(2p^\mu)u$ at $p'=p$ & Prob.~38.3 & $<10^{-14}$ \\
  $\{\gamma^5,\gamma^\mu\}=0$ for all $\mu$ & (38.41 prereq.) & $<10^{-14}$ \\
  $\sum_s u_s\bar u_s = -\slashed{p}+m$ & (38.28) & $<10^{-14}$ \\
  $\sum_s v_s\bar v_s = -\slashed{p}-m$ & (38.29) & $<10^{-14}$ \\
  $\bar u_{s'}\gamma^\mu u_s = 0$ for $s\neq s'$ & (38.21) & $<10^{-14}$ \\
  $\bar u_{s'}(p)\gamma^0 v_s(-p)=0$ & (38.22) & $<10^{-14}$ \\
  $P_{L,R}^2=P_{L,R}$, $P_LP_R=0$ & chirality projectors & $<10^{-14}$ \\
  $P_Lu_- = u_-$, $P_Ru_+ = u_+$ & massless chirality & $<10^{-14}$ \\
  \bottomrule
\end{tabular}
\end{center}

All identities verified to machine precision ($\lesssim 10^{-15}$).

%% ============================================================
\section{Summary}
%% ============================================================

\begin{center}
\renewcommand{\arraystretch}{1.5}
\begin{tabular}{@{}ll@{}}
  \toprule
  Object & Expression \\
  \midrule
  $\sigma^\mu$ completeness & $\sigma^\mu_{\alpha\dot\alpha}\sigma_\mu^{\beta\dot\beta}=-2\delta_\alpha^\beta\delta_{\dot\alpha}^{\dot\beta}$ \\
  Alternative form & $\varepsilon^{\alpha\beta}\varepsilon^{\dot\alpha\dot\beta}\sigma^\mu_{\alpha\dot\alpha}\sigma^\nu_{\beta\dot\beta}=-2g^{\mu\nu}$ \\
  2-trace & $\mathrm{Tr}[\sigma^\mu\bar\sigma^\nu]=\mathrm{Tr}[\bar\sigma^\mu\sigma^\nu]=-2g^{\mu\nu}$ \\
  4-trace & $\mathrm{Tr}[\sigma^\mu\bar\sigma^\nu\sigma^\rho\bar\sigma^\sigma]=2(g^{\mu\nu}g^{\rho\sigma}-g^{\mu\rho}g^{\nu\sigma}+g^{\mu\sigma}g^{\nu\rho})+2i\varepsilon^{\mu\nu\rho\sigma}$ \\
  Index lowering & $\psi_\alpha=\varepsilon_{\alpha\beta}\psi^\beta$,\quad $\varepsilon_{12}=-1$ \\
  Index raising & $\psi^\alpha=\varepsilon^{\alpha\beta}\psi_\beta$,\quad $\varepsilon^{12}=+1$ \\
  $\sigma\leftrightarrow\bar\sigma$ & $\bar\sigma^{\mu\dot\alpha\alpha}=\varepsilon^{\alpha\beta}\varepsilon^{\dot\alpha\dot\beta}\sigma^\mu_{\beta\dot\beta}$ \\
  Schouten & $\varepsilon_{\alpha\gamma}\varepsilon_{\beta\delta}=\varepsilon_{\alpha\beta}\varepsilon_{\gamma\delta}-\varepsilon_{\alpha\delta}\varepsilon_{\gamma\beta}$ \\
  $\sigma$-Fierz & $\bar\sigma^{\mu\dot\alpha\alpha}\sigma_{\mu\beta\dot\beta}=-2\delta^\alpha_\beta\delta^{\dot\alpha}_{\dot\beta}$ \\
  4-Fermi Fierz & $(\chi^\dagger\bar\sigma^\mu\psi)(\xi^\dagger\bar\sigma_\mu\eta)=-2(\chi^\dagger\xi^\dagger)(\eta\psi)$ \\
  Gordon identity & $2m\bar u'\gamma^\mu u = \bar u'[(p'+p)^\mu - 2iS^{\mu\nu}(p'-p)_\nu]u$ \\
  Extended Gordon & $\bar u'[\cdots]\gamma^5 u = 0$;\quad $\{\gamma^5,\gamma^\mu\}=0$ \\
  Spin-sum (particle) & $\sum_s u_s\bar u_s = -\slashed{p}+m$ \\
  Spin-sum (antiparticle) & $\sum_s v_s\bar v_s = -\slashed{p}-m$ \\
  Massless helicity & $u_\pm\bar u_\pm = P_{R/L}(-\slashed{p})$;\quad helicity = chirality \\
  \bottomrule
\end{tabular}
\end{center}

\bigskip
\noindent\textbf{Next:} Chapter 48 (spinors for massless particles) builds
on these tools toward the spinor-helicity formalism and MHV amplitudes.
